\chapter*{Conclusion and Future Research}
\addcontentsline{toc}{chapter}{Conclusion and Future Research}
\markboth{Conclusion and Future Research}{Conclusion and Future Research}
\chapterstatus{Unnumbered synthesis of findings, present limitations and
genuinely new research directions. Items already completed by the unified
project are recorded as results rather than repeated as future work.}

\section*{What the project has completed}

This thesis converts eight specialised studies into one traceable research
architecture. Its main contribution is the separation and subsequent
connection of five objects: the historical gold series, the option-implied
gold-price law, production, unit cost and corporate cash flow. The resulting
chain is explicit enough to identify which evidence supports each output and
where a change of data, units, measure or accounting perimeter occurs.

Several tasks described as future work in the standalone articles are now
complete. The Team 8 project no longer stops at Black--Scholes, Heston or
independent Bates jumps: it contains the full affine Bates--Hawkes
specification, constrained calibration, residual and smile diagnostics,
four-model path simulation and a dense-only rolling comparison over 30 origins and 4,667 common forecasts. Its
price scenarios have also been connected to Barrick 2026Q1--Q2 actuals, the
Team~4 forecast from 2026Q3 and the unified DCF/WACC/equity contract. The
Team~4 price demonstrator is excluded. The valuation is therefore an implemented
four-branch experiment, not a proposed future integration.

Under that common contract, the four conditional Barrick medians lie between
USD~35.14 and USD~35.87 per share, below the completed USD~44.13 close
on 2 September 2026. The close remains inside every broad distribution,
with 40.76--41.92\% of simulated values above it. Full Bates--Hawkes is the
primary OOS-ranked structural scenario and gives a median of USD~35.14,
whereas Heston is the current-snapshot fit winner. The shared direction
is a property of the admitted assumptions, not evidence of market mispricing.

The comparison further shows that changing the gold-price law affects the
tails more than the central corporate value when production, cost, WACC and
the equity bridge are fixed. This is a useful attribution result: it identifies
where price-model complexity propagates through the DCF and where downstream
assumptions compress or dominate its effect.

\section*{Present limitations}

The completed integration remains conditional at several economic interfaces.
First, GLD option dynamics are estimated under \(\mathbb Q\), whereas a
long-horizon operating valuation would ideally use scenarios justified under
\(\mathbb P\). Restarting the option-implied engines at a physical-gold level
transfers distributional shape; it is not a validated GLD-to-gold conversion
or change of measure. The gold start, Treasury curve, calibration surface,
Barrick close and current surface audit are explicitly dated: the market snapshot is 2 September, while the realised-price anchor is a Q2 average.

Second, the unified run uses Barrick actuals only for 2026Q1--Q2 and fixes the
Team~4 forecast thereafter. That design isolates price-engine risk but does not propagate reserve,
grade, recovery, energy, labour, mine-mix or closure uncertainty. Dependence
between gold prices and operating conditions is therefore not estimated.

Third, Team 5 supplies a coherent DCF architecture rather than a
filing-reconciled Barrick model. The operating-to-EBIT bridge, cash taxes,
working capital, sustaining and growth capex, debt, cash, minorities, streams,
other claims, non-operating assets and diluted shares require one dated
corporate reconciliation. The present zero bridge items and share-count proxy
are assumptions. Finite-mine closure also remains less developed than the
generic going-concern terminal block.

Finally, a sophisticated stochastic law does not eliminate model risk. The
current cross-section favours Heston, whereas Full Bates--Hawkes leads the rolling OOS ranking. Its branching ratio is 0.5569, not the former August floor. No model beats interpolated IV persistence. Parameter pressure,
surface sampling, optimizer dependence, regime shifts and the choice of
physical-measure transformation remain sources of uncertainty.

Costs retain a persistent trend without additional technology-driven reductions.
This is prudential relative to otherwise identical efficiency-improvement
scenarios, not a mathematical lower bound. Stochastic operating costs would
introduce favourable and adverse outcomes, and their conditional mean could
remain smooth. Productivity improvements and their investment costs must be
evaluated jointly with the growth already included in the DCF.

\section*{Future research programme}

The next research cycle should begin with economic reconciliation rather than
another layer of stochastic complexity. A dated Barrick data contract should
map mine-level production, realised prices, royalties, cash costs, AISC,
depreciation, sustaining and growth capex, working capital, taxes and closure
liabilities into EBIT, invested capital and FCFF. Debt, cash, minority
interests, streams, non-operating assets and diluted shares should then close
the EV-to-equity bridge to the same filing and listing.

The commodity interface requires a separate empirical programme. The GLD-to-
physical-gold basis should be estimated and stress-tested, and risk-neutral
option scenarios should be converted into physical scenarios through an
explicitly validated risk-premium or filtering model. Only after that step can
the long-horizon paths be interpreted as more than option-implied conditional
shapes.

Operating uncertainty should then be reintroduced jointly. Team 4 production
and cost models can be sampled alongside price, but their dependence structure
must be estimated or scenario-defined. Mine depletion, expansion projects,
reserve conversion, inflation, energy and foreign-exchange shocks should feed
an explicit finite-resource closure. The comparison can then be repeated to
measure whether gold-model choice or operating/accounting uncertainty explains
more of the equity-value distribution.

A later information layer may augment the filtration with timestamped
macroeconomic and geopolitical data. News embeddings, NLP classifications or
machine-learning forecasts are admissible only with event-time alignment,
leakage controls, pre-declared benchmarks and true chronological evaluation.
The relevant question is not whether a flexible model improves in-sample fit,
but whether incremental information improves forecasts beyond price,
volatility and persistence benchmarks.

Further model research may examine time-varying risk premia, regime-dependent
parameters, model averaging and a rough-Hawkes option-pricing construction.
Those extensions should be judged through the same hierarchy used here:
mathematical admissibility, deterministic implementation tests, dated
out-of-sample evidence and economic propagation through a reconciled corporate
model.

\section*{Concluding synthesis}

The project demonstrates the value of interdisciplinarity in applied finance.
Probability, stochastic calculus, econometrics, numerical methods,
optimization, accounting and corporate valuation are not parallel subjects in
this setting; each constrains a different interface of the same experiment.
The strongest outcome is therefore not one point estimate or one preferred
model. It is a reproducible framework that keeps historical evidence,
option-implied information, operating assumptions and accounting claims
separate until the moment they are deliberately connected.

Within that framework, Full Bates--Hawkes offers the richest admitted
description of current option-implied variance and clustered jumps, while the
four-model valuation shows how little a better surface fit alone can establish
about a mining company's equity. The thesis closes with a conditional result,
an explicit inventory of unresolved economic bridges and a research programme
capable of testing them. It does not issue a fair-value certification, target
price or investment recommendation.
